\documentclass[11pt]{article}

\usepackage[margin=1.6cm]{geometry}
\usepackage{amsmath}
\usepackage{amssymb}
\usepackage{enumitem}
\usepackage{titlesec}
\usepackage{xcolor}
\usepackage{fancyhdr}
\usepackage{multicol}
\usepackage{array}

\titleformat{\section}{\large\bfseries\color{blue!50!black}}{\thesection}{0.6em}{}
\titlespacing{\section}{0pt}{1.1em}{0.4em}
\titleformat{\subsection}{\normalsize\bfseries}{\thesubsection}{0.6em}{}
\titlespacing{\subsection}{0pt}{0.8em}{0.3em}

\setlength{\headheight}{14pt}

\pagestyle{fancy}
\fancyhf{}
\fancyhead[L]{Year 9 Revision: FDP and Statistics}
\fancyfoot[C]{\thepage}
\renewcommand{\headrulewidth}{0.4pt}

\fancypagestyle{firstpage}{%
  \fancyhf{}
  \fancyhead[L]{Year 9 Revision: FDP and Statistics}
  \fancyhead[R]{Name: \makebox[4cm]{\hrulefill}}
  \fancyfoot[C]{\thepage}
  \renewcommand{\headrulewidth}{0.4pt}
}

\newcommand{\ans}{\hrulefill}
\newcommand{\anshere}{\par\vspace{0.4em}\noindent\ans\par\vspace{0.4em}}
\setlist[enumerate,1]{leftmargin=1.4em, itemsep=0.5em, topsep=0.3em}
\setlength{\columnsep}{1.1cm}
\setlength{\columnseprule}{0.4pt}
\raggedcolumns

\newcolumntype{Y}{>{\centering\arraybackslash}p{2.3cm}}

\begin{document}
\thispagestyle{firstpage}

\begin{center}
  {\Large\bfseries Year 9 Revision: FDP and Statistics}
\end{center}
\vspace{0.3em}

\section{Fractions, Decimals and Percentages}

\noindent\textbf{Converting between forms:}
\begin{itemize}[leftmargin=1.4em, itemsep=0.2em, topsep=0.2em]
  \item Fraction $\to$ decimal: divide the numerator by the denominator.
  \item Decimal $\to$ percentage: multiply by $100$ (move the decimal point $2$ places right).
  \item Percentage $\to$ fraction: write over $100$, then simplify.
\end{itemize}

\vspace{0.3em}
\noindent\textbf{Worked example:} Convert $\dfrac{3}{4}$ to a decimal and a percentage.
\[
  \frac{3}{4} = 3 \div 4 = 0.75, \qquad 0.75 \times 100 = 75\%.
\]

\vspace{0.4em}
\noindent Complete the table below, giving fractions in simplest form.

\begin{center}
\begin{tabular}{|Y|Y|Y|}
\hline
\textbf{Fraction} & \textbf{Decimal} & \textbf{Percentage} \\\hline
$\dfrac{3}{4}$ & $0.75$ & $75\%$ \\[0.6em]\hline
$\dfrac{2}{5}$ & & \\[0.9em]\hline
 & $0.7$ & \\[0.9em]\hline
$\dfrac{1}{8}$ & & \\[0.9em]\hline
 & $0.45$ & \\[0.9em]\hline
 & & $37.5\%$ \\[0.9em]\hline
$\dfrac{11}{25}$ & & \\[0.9em]\hline
 & & $34\%$ \\[0.9em]\hline
\end{tabular}
\end{center}

\newpage

\subsection*{Using FDP in probability}
\noindent Probability can be written as a fraction, decimal, \emph{or} percentage --- they all
mean the same thing.
\[
  P(\text{event}) = \frac{\text{favourable outcomes}}{\text{total outcomes}}.
\]

\vspace{0.3em}
\noindent\textbf{Worked example:} A bag has $5$ red, $3$ blue and $2$ green marbles ($10$ in
total). Find $P(\text{red})$ as a fraction, decimal and percentage.
\[
  P(\text{red}) = \frac{5}{10} = \frac{1}{2} = 0.5 = 50\%.
\]

\begin{multicols}{2}
\begin{enumerate}[label=(\alph*)]
  \item Using the same bag of marbles, find $P(\text{blue})$ as a fraction, decimal and
    percentage.
  \anshere
  \item Find $P(\text{green})$ as a fraction, decimal and percentage.
  \anshere
  \item Find $P(\text{not red})$ as a fraction, decimal and percentage. (Hint: use your answer
    to the worked example.)
  \anshere
  \item A spinner has $8$ equal sections: $3$ blue, $4$ yellow, $1$ red. Find $P(\text{blue})$
    as a fraction, decimal and percentage.
  \anshere
  \item Find $P(\text{yellow})$ for the same spinner.
  \anshere
  \item A standard die is rolled. Find $P(\text{even number})$ as a fraction, decimal and
    percentage.
  \anshere
  \item Find $P(\text{a number more than } 4)$ for the same die. Round your decimal and
    percentage to $1$ decimal place.
  \anshere
  \item A card is drawn from a standard deck of $52$ cards. Find $P(\text{a heart})$ as a
    fraction, decimal and percentage. (There are $13$ hearts in a deck.)
  \anshere
\end{enumerate}
\end{multicols}

\vspace{0.3em}
\begin{enumerate}[label=(\alph*), resume]
  \item A vending machine has a $15\%$ chance of jamming each time it is used. Find the
    probability that it does \emph{not} jam, as a percentage, a decimal, and a fraction in
    simplest form.
  \anshere
  \item A weather forecast says there is a ``$30\%$ chance of rain'' tomorrow. Write this
    probability as a fraction and a decimal. Then explain in a sentence what a $30\%$ chance of
    rain actually means --- does it guarantee it will \emph{not} rain?
  \anshere
\end{enumerate}

\newpage

\section{Statistics Revision}

\subsection*{Sound statistical questions}
\noindent A \textbf{statistical question} is one that can be answered by \emph{collecting data},
where you expect the data to \textbf{vary}. It usually asks about a \emph{group}, not a single
fixed fact.
\[
  \text{\emph{Not statistical:}} \;\; \text{``How old am I?''} \qquad
  \text{\emph{Statistical:}} \;\; \text{``What are the ages of people in my family?''}
\]
The first question has one fixed answer. The second expects a range of different answers, which
you would need to collect data to find.

\vspace{0.4em}
\noindent\textbf{For each question below, write S (statistical) or N (not statistical), and
give a one-sentence reason.}

\begin{multicols}{2}
\begin{enumerate}[label=(\alph*)]
  \item How many siblings does Jamie have?
  \anshere
  \item How many siblings do students in our class have?
  \anshere
  \item What is the temperature outside right now?
  \anshere
  \item What were the daily high temperatures in Wellington this month?
  \anshere
  \item What is the most popular pizza topping among Year 9 students at our school?
  \anshere
  \item What is my favourite pizza topping?
  \anshere
  \item How long is the Auckland Harbour Bridge?
  \anshere
  \item How long does it take different students to run $100$\,m?
  \anshere
\end{enumerate}
\end{multicols}

\vspace{0.3em}
\noindent\textbf{Rewrite each question below so that it becomes a sound statistical question.}
\begin{enumerate}[label=(\alph*)]
  \item ``How tall am I?''
  \anshere
  \item ``What is my favourite sport?''
  \anshere
  \item ``How many hours of sleep did I get last night?''
  \anshere
\end{enumerate}

\newpage

\subsection*{Explanatory vs response variables}
\noindent When we investigate whether two variables are related, the \textbf{explanatory
variable} is the one we think might explain or influence a change. The \textbf{response
variable} is the one being measured --- it may \emph{respond} to changes in the explanatory
variable. (This does not always mean one \emph{causes} the other --- only that we are
investigating a possible relationship.)

\vspace{0.3em}
\noindent\textbf{Worked example:} A scientist wants to find out if the amount of sunlight a
plant gets affects how tall it grows.
\begin{itemize}[leftmargin=1.4em, itemsep=0.2em, topsep=0.2em]
  \item Explanatory variable: amount of sunlight.
  \item Response variable: height of the plant.
\end{itemize}

\vspace{0.4em}
\noindent\textbf{For each scenario, state the explanatory variable and the response variable.}
\begin{enumerate}[label=(\alph*)]
  \item Does the number of hours spent studying affect exam scores?
  \anshere
  \item Is there a relationship between a person's age and their reaction time?
  \anshere
  \item Does the amount of water in a kettle affect how long it takes to boil?
  \anshere
  \item Researchers want to know if the amount of fertiliser given to tomato plants affects the
    number of tomatoes produced.
  \anshere
  \item Does the price of a movie ticket affect how many people attend the cinema?
  \anshere
  \item Is a student's typing speed related to how many hours per week they spend using a
    computer?
  \anshere
\end{enumerate}

\vspace{0.4em}
\noindent\textbf{Putting it together:} Write your own statistical question that investigates a
relationship between two variables. Identify the explanatory variable and the response variable
in your question.
\anshere
\anshere

\end{document}
